% id: 0473
% book: RPM 확률과 통계 · 05 확률분포 (2)
% section: 유형 11 이항분포와 정규분포의 관계의 활용; 확률 구하기
% figure: crop:fig-0473.png
% answer: ②
% answer_source: 답지(쪽 렌더)
% variant_level: 0 (원본 전사)
% 이 파일은 build.mjs 가 items.json 에서 만든다. 고칠 때는 items.json 을 고친다.
\smash{\raisebox{1.5mm}{\dmkichul{RPM 확률과 통계 0473번 · 78쪽}}}
\noindent\begin{minipage}[t]{0.58\linewidth}\vspace{0pt}\raggedright 흰 공 $1$개와 검은 공 $3$개가 들어 있는 주머니에서 한 개의 공을 꺼내어 색을 확인하고 다시 넣는 시행을 $192$번 반복할 때, 흰 공이 $57$번 이상 나올 확률을 오른쪽 표준정규분포표를 이용하여 구한 것은?\end{minipage}\hfill\begin{minipage}[t]{0.4\linewidth}\vspace{0pt}\centering \setlength{\dmfigw}{106.7pt}\ifdim\dmfigw>\linewidth\setlength{\dmfigw}{\linewidth}\fi\includegraphics[width=\dmfigw]{figures/fig-0473.png}\end{minipage}\par
\begin{choices32}{$0.0228$}{$0.0668$}{$0.1587$}{$0.8413$}{$0.9332$}\end{choices32}
